% Generated draft: requires mathematical and source review.
\documentclass[11pt]{article}
\usepackage[margin=25mm]{geometry}
\usepackage{fontspec}
\setmainfont{Latin Modern Roman}
\usepackage{amsmath,amssymb,mathtools,unicode-math}
\setmathfont{Latin Modern Math}
\usepackage{enumitem,needspace,xurl,hyperref}
\setlength{\parindent}{0pt}
\setlength{\parskip}{5pt}
\setlength{\emergencystretch}{3em}
\newcommand{\N}{\mathbb{N}}
\newcommand{\Z}{\mathbb{Z}}
\newcommand{\Q}{\mathbb{Q}}
\newcommand{\R}{\mathbb{R}}
\newcommand{\C}{\mathbb{C}}
\newcommand{\F}{\mathbb{F}}
\newcommand{\A}{\mathbb{A}}
\newcommand{\PP}{\mathbb{P}}
\newcommand{\E}{\mathbb{E}}
\newcommand{\eps}{\varepsilon}
\newcommand{\dd}{\,\mathrm{d}}
\begin{document}
\section*{UnsolvedMath ID: 5000005}
\textbf{Source problem code:} catalogue entry
\subsection{Definitions and mathematical statement}
Let $n\ge3$ be an integer and consider billiards in a regular $n$-gon. Unfold a billiard path by reflecting the polygon across each side it encounters. A short trajectory is a segment joining vertices in this unfolding and meeting no vertex in its interior; a reachable point is an endpoint vertex that occurs in such an unfolding.

For an oriented short trajectory, the source associates source and target endpoint angles, denoted $\alpha$ and $\beta$ in its Definition 2.1, and one of the types
\[
 A_k,\qquad k\in\Z/(n-2)\Z.
\]
The supplied record gives only the particular characterization
\[
 A_0:\qquad \beta-\alpha=\frac{2\pi}{n}.
\]
It does not supply Definition 2.1, the other type criteria, or the coordinate convention for the distinguished point $O$. These are essential dependencies of the proposed type rule.

Write $\theta$ for the slope angle called $\alpha$ in the transformation clause. Conjecture 2.6 asks that a short trajectory from $O$ of type $A_k$ and slope angle $\theta$ transform under
\[
 \theta\longmapsto\theta'=\frac{\ell\pi}{n}+\varepsilon\theta,
 \qquad\varepsilon\in\{-1,1\},
\]
according to
\[
 A_k\longmapsto A_{\varepsilon k-\ell},
\]
where type indices are reduced modulo $n-2$. For odd $n$, $\ell$ ranges over integers; for even $n$, $\ell$ must be even. The assertion concerns the short trajectory from $O$ at the transformed angle, including the presumed existence of that trajectory in the intended construction, not just a formal transformation of indices.

The source calls that transformed trajectory ``parallel'' although the displayed angles need not be parallel in the usual Euclidean sense. It also reuses $\alpha$ for an endpoint angle and a slope angle without explaining the identification. The relevant trajectory-family meaning and angle conventions must be recovered from the source; neither ordinary parallelism nor an equality of differently defined angles is silently imposed.
\subsection{Short English statement}
The proposed symmetry rule sends type $A_k$ to $A_{\varepsilon k-\ell}$ when the slope changes to $\ell\pi/n+\varepsilon\theta$, with even $\ell$ required for even $n$. The missing endpoint-angle and type definitions prevent a complete standalone geometric formulation.

The supplied summary places this in Fuchs's conjectural family 2.3--2.7 and reports no later resolution found through 2026. The registry labels it open with OPEN-TRIAGE; provenance records no later authoritative resolution check. Those assessments are retained without an independent source visit or status claim.
\subsection{Sources}
\begin{itemize}
\item Local UnsolvedMath registry, permanent ID 5000005, source code catalogue entry, ``Types of parallel short trajectories,'' proposed by Dmitry Fuchs, 2020. Dataset curation and metadata: Ulam AI and UnsolvedMath Contributors, CC BY 4.0.
\item Dmitry Fuchs, \emph{Billiard Trajectories in Regular Polygons and Foliations by Closed Geodesics on Surfaces} (2020), Conjecture 2.6 and referenced Definition 2.1, \url{https://amj.math.stonybrook.edu/pdf-Springer-final/020-0170.pdf}; pdf-text extraction, accessed 2026-07-29.
\item The supplied research summary places the conjecture in the family 2.3--2.7 but gives no later primary resolution reference.
\item Source provenance marks status evidence NEEDS\_REVIEW and rights NEEDS\_REVIEW: the source presents the item as open, no later authoritative resolution check is recorded, and redistribution permission was not established despite verified public access. This qualification is separate from the dataset licence.
\end{itemize}
\end{document}
